\movedsub{The Correctness Grader Admits Refusals to the Wrong Population}
\label{sec:refusals}

The same problem returns once faithfulness is conditioned on correctness. String containment marks an answer correct when a reference string occurs in it, and a short reference string can occur inside a refusal; on a 200-answer calibration ten refusals were marked correct this way and the judge overturned nine. Each such row carries a near-floor score into the correct-answer average and lowers the baseline every correct-versus-incorrect comparison is measured against. Under graded judgement no refusal is correct in any of the twenty-four judged cells.

That is enough to change a sign. With containment, and a baseline excluding refusals from the incorrect side only, NQ/Claude answers looked more grounded when wrong than when right by 0.018 to 0.050. Excluding refusals from both sides gives -0.006 to +0.018; grading with the judge gives -0.051 to -0.113. Each step corrects a grader, not an answer, and all three move the quantity the same way. Tested directly over all twenty-four judged cells, a permutation test with Holm correction finds six significant, every one of them negative, and none of the four positive-signed cells has an interval excluding zero. E5-instruct on Natural Questions is negative under all three generators (-0.105, -0.115, -0.125) and survives correction in each. We report no evidence that answers are more grounded when wrong. Intervals are 0.064 to 0.129 wide on 44 to 134 wrong answered rows per cell, so this bears on the general claim rather than on individual systems.

The same accounting applies to the contingency between retrieval success and correctness (Appendix~\ref{app:cell}). The cell where retrieval succeeded and the answer was still wrong bears on whether retrieval quality is sufficient, and 69 to 87 per cent of it was never attempted. Read whole it runs from 17 to 36 per cent of queries and would support the claim that good retrieval often fails to produce a correct answer. Most of what it counts is a system declining to answer, and refusal rates differ by more than twenty points across the systems compared. What is left is real, and 4.5 to 7.9 per cent.
